\documentclass[english,twoside, openright, hidelinks, 11 pt, class=report,crop=false]{standalone}
\input{../../preamb}
\input{../bm}

\begin{document}	
\section{Definisjoner}\index{vektor}
En \outl{todimensjonal vektor} angir en forflytning i et koordinatsystem med en $ x $-akse og en $ y $-akse. En vektor tegner vi som et linjestykke mellom to punkt, i tillegg til at vi lar en pil vise til hva som er endepunktet. Det betyr at forflytningen starter i punktet uten pil, og ender i punktet med pil.
\begin{figure}
	\centering
\subfloat[]{\includegraphics{\figp{vekdefa}}}\qquad
\subfloat[]{\includegraphics{\figp{vekdefb}}}
\end{figure}
I figur \textsl{(a)} er vektoren $ \vec{v} $ vist med startpunkt $ (0, 0) $ og endepunkt $ (2,3) $. Når en vektor har startpunkt $ (0, 0) $, sier vi at den er vist i \outl{grunn-stillingen}. I figur \textsl{(b)} er $ \vec{v} $ vist med startpunkt $ (1, -2) $ og endepunkt $ (3, 1) $. Forflytningen $ \vec{v} $ viser til er å vandre 2 mot høyre langs $ x $-aksen og $ 3 $ opp langs $ y $-aksen. Dette skriver vi som $ {\vec{v}=[2, 3]} $, som er $ \vec{v} $ skrevet på \outl{komponentform}. 2 og 3 er henholdvis $ x $-komponenten og $ y $-komponenten til $ \vec{v} $.\vsk

\spr{En vektor med to komponenter kalles en \outl{todimensjonal \\ vektor}  eller en \outl{vektor i planet}. \vsk

\textit{Komponent} kan i mange sammenhenger ses på som et synonym for \textit{koordinat}.
}
\newpage
\eks[1]{
\[ \vec{a}= [1, 3]\quad,\quad \vec{b} = [0, -2] \quad,\quad
\vec{c}= [-3, -4] \quad,\quad \vec{d}= [5, 0] \]		
\fig{vek1}
}
\newpage
\regdef[Vektoren mellom to punkt]{
En vektor $ \vec{v} $ med startpunkt $ {(x_1, y_1)} $ og endepunkt $ (x_2, y_2) $ er gitt som 
\begin{equation*}
	\vec{v}= [x_2-x_1, y_2-y_1]
\end{equation*}
}
\eks[1]{
Skriv vektorene på komponentform.\os

$ \vec{a} $ har startpunkt $ (1, 3) $ og endepunkt $ (7, 5) $. \os
$ \vec{b} $ har startpunkt $ (0, 9) $ og endepunkt $ (-3, 2) $. \os
$ \vec{c} $ har startpunkt $ (-3, 7) $ og endepunkt $ (2, -4) $. \os
$ \vec{d} $ har startpunkt $ (-7, -5) $ og endepunkt $ (3, 0) $.

\sv \vs

\algv{
\vec{a}&=[7-1, 5-3]=[6, 2] \br
\vec{b}&=[-3-0, 2-9]=[-3, -7] \br
\vec{c}&=[2-(-3), -4-7]=[5, -11]\br
\vec{d}&=[3-(-7),0-(-5)]=[10, 5]
}
}
\info{Punkt eller vektor}{
Rent matematisk er det ingen forskjell på et punkt og en vektor; punktet $ (a, b) $ viser til akkurat samme plassering som vektoren $ [a, b] $, og begge kan vise til den samme forflytningen. Ofte kan det likevel være greit å skille mellom når vi snakker om en plassering og når vi snakker om en forflytning, og til det bruker vi begrepene punkt (plassering) og vektorer (forflytning).
}
\section{Regneregler for vektorer}
\regdef[Addisjon og subtraksjon av vektorer]{
	Gitt vektorene $ {\vec{u}=[x_1, y_1]} $ og $ {\vec{v}=[x_2, y_2]} $, og punktet ${ A=(x_0, y_0)} $. Da er 
	\begin{align*}
		A+\vec{u} &= (x_0+x_1, y_0+y_1)	\br
		\vec{u}+\vec{v} &= [x_1+x_2, y_1+y_2] \br
		\vec{u}-\vec{v} &= [x_1-x_2, y_1-y_2] 
	\end{align*}
	Summen eller differansen av $ \vec{u} $ og $ \vec{v} $ kan vi tegne slik:
	\begin{figure}
		\centering
		\includegraphics[scale=1]{\figp{uplusv}}\qquad\quad
		\includegraphics[scale=1]{\figp{uminusv}}
	\end{figure}
}
\reg[Regneregler for vektorer]{
For vektorene $\vec{u}=[x_1, y_1]$, $ \vec{v} $ og $\vec{w} $, og et tall $ t $, har vi at
\begin{align*}
	t\vec{u} &= [tx_1, ty_1] \br
	t(\vec{u}+\vec{v})&= t\vec{u}+t\vec{v}\br
	(\vec{u}+\vec{v})+\vec{w} &= \vec{u}+(\vec{v}+\vec{w}) \br
	\vec{u}-(\vec{v}+\vec{w}) &= \vec{u}-\vec{v}-\vec{w}
\end{align*}
}
\eks[1]{
Skriv uttrykket så enkelt som mulig.
\[ t(\vec{u}+\vec{v})-(t\vec{u}+\vec{v}) \] \vs \vs

\sv \vsb
\begin{align*}
	t(\vec{u}+\vec{v})-(t\vec{u}+\vec{v}) = t\vec{u}+t\vec{v}-t\vec{u}-\vec{v}=(t-1)\vec{v}
\end{align*}
}

\newpage
\eks[2]{
I trekanten under er $AB=5AD$ og $BC=4BF$. \\ Finn $\vv{AC}$ uttrykt ved $\vec{u}$ og $\vec{w}$.
\fig{opgvekuvabcexmp}
\sv
Vi har at
\begin{align*}
	\vv{AC}&=\vv{AB}+\vv{BC} \\
	&=5\vv{AD}+4\vv{BF} \\
	&=5\vec{u}+4\cdot(-\vec{v})\\
	&=5\vec{u}-4\vec{v}
\end{align*}
}
\regdef[Vinkelen mellom to vektorer]{
Vinkelen mellom to vektorer er (den minste) vinkelen som blir dannet når vektorene plasseres i samme startpunkt. For to vektorer $ \vec{u} $ og $ \vec{v} $ skriver vi denne vinkelen som $ \angle(\vec{u}, \vec{v}) $.
\fig{vekvink}}
\info{Vinkelmål}{I vektorregning er det vanlig å oppgi vinkler i grader, altså på\\ intervallet $ [0^\circ, 180^\circ] $.}
\section{Lengden til en vektor}
Gitt en vektor $ \vec{v}=[x_1, y_1] $. \outl{Lengden} til $ \vec{v} $ er avstanden mellom startpunktet og endepunktet. 
\begin{figure}
\centering
\subfloat[]{\includegraphics{\figp{veklena}}}\qquad\quad
\subfloat[]{\includegraphics{\figp{veklenb}}}
\end{figure}
Av envektor kan vi danne en rettvinklet trekant hvor $ |\vec{v}| $ er lengden til hypotenusen, og $ |x_1| $ og $ |y_1| $ er de respektive lengdene til katetene. Dermed er $ |\vec{v}| $ gitt av Pytagoras' setning.\regv

\reg[Lengden til en vektor \label{veklength}]{
Gitt en vektor $ \vec{v}=[x_1, y_1] $. Lengden $ |\vec{v}| $ er da 
\begin{equation*}
	|\vec{v}|=\sqrt{x_1^2+y_1^2}
\end{equation*}
}
\eks[1]{
Finn lengdene til vektorene $ \vec{a}= [7, 4] $ og $ \vec{b}= [-3, 2] $.

\sv 
\algv{
	|\vec{a}|&=\sqrt{7^2+4^2}=\sqrt{65} \br
	|\vec{b}|&=\sqrt{(-3)^2+2^2}=\sqrt{13}
}
}

\section{Skalarproduktet}
\subsection*{Skalarproduktet I}
\regdef[Skalarproduktet I \label{scalardef}]{
For to vektorer $ {\vec{u}=[x_1, y_1]} $ og $ {\vec{v}=[x_2, y_2]} $, er \outl{skalarproduktet} gitt som
\begin{equation*}
	\vec{u}\cdot\vec{v}=x_1x_2+y_1y_2
\end{equation*}
}
\spr{
Skalarproduktet kalles også \outl{prikkproduktet} eller \\ \outl{indreproduktet}.\vsk

Ordet \textit{skalar} viser til en éndimensjonal størrelse, altså et tall.
}
\eks[1]{
Gitt vektorene $ \vec{a}=[3, 2] $, $ \vec{b}=[4,7] $ og $ \vec{c}=[1,-9] $. Regn ut $ \vec{a}\cdot\vec{b} $ og $ \vec{a}\cdot\vec{c} $.

\sv 
\algv{
\vec{a}\cdot\vec{b}&=3\cdot4+2\cdot7=26 \br
\vec{a}\cdot\vec{c}&=3\cdot1+2(-9)=-15
}
}
\newpage
\reg[Regneregler for skalarproduktet \label{regnerules}]{
	For vektorene $ \vec{u} $, $ \vec{v} $ og $ \vec{w} $ har vi at
\begin{align*}
	\vec{u}\cdot \vec{u}&=\vec{u}^{\,2} \label{veku2}\br
	\vec{u}\cdot\vec{v} &= \vec{v}\cdot\vec{u} \br
	\vec{u}\cdot\left(\vec{v}+\vec{w}\right) &= \vec{u}\cdot\vec{v}+\vec{u}\cdot\vec{w} \br
	\left(\vec{u}+\vec{v}\right)^2 &= \vec{u}^{\,2} + 2\vec{u}\cdot\vec{v}+\vec{v}^{\,2}
\end{align*}
}
\eks[]{Forkort uttrykket
	\[   \vec{b}\cdot(\vec{a}+\vec{c}) + \vec{a}\cdot(\vec{a}+\vec{b})+\vec{b}^{\,2} \]
	
	når du vet at $ \vec{b}\cdot\vec{c}=0 $. \\
	
	\sv
	\algv{\vec{b}\cdot(\vec{a}+\vec{c}) + \vec{a}\cdot(\vec{a}+\vec{b})+\vec{b}^{\,2} &= \vec{b}\cdot\vec{a}+\vec{b}\cdot\vec{c} + \vec{a}\cdot\vec{a}+\vec{a}\cdot\vec{b}+\vec{b}^2 \\
		&= \vec{a}^{\,2} + 2 \vec{a}\cdot \vec{b} + \vec{b}^{\,2} \\
		&= \left(\vec{a}+ \vec{b}\right)^2
	}
} 

\subsection*{Skalarproduktet II}
Gitt vektoren $ \vec{u}-\vec{v} $, hvor 
$ \vec{u}=[x_1, y_1] $ og $ \vec{v}=[x_2, y_2] $. Da er
\[ \vec{u}-\vec{v}=[x_1-x_2, y_1-y_2] \]
Av \rref{veklength} har vi at
\begin{align}
	|\vec{u}-\vec{v}|&= \sqrt{(x_1-x_2)^2 + (y_1-y_2)^2} \nonumber \\
	&= \sqrt{x_1^2 - 2x_1 x_2 + x_2^2 + y_1^2 - 2y_1 y_2 + y_2^2 }\label{preskal}
\end{align}

Ved hjelp av \rref{scalardef} og \rref{regnerules} kan vi skrive \eqref{preskal} som
\begin{equation}
	|\vec{u}-\vec{v}|= \sqrt{\vec{u}^{\,2} - 2\vec{u}\cdot\vec{v} + \vec{v}^{\,2}} \label{skl1}
\end{equation}
Videre merker vi oss følgende figur:
\fig{skaldef}
Av \refunnbr{cossetn}{cosinussetningen} (\rref{cossetn}) og \eqref{skl1} er
\alg{
	|(\vec{v}-\vec{u})|^2&= |\vec{v}|^2+|\vec{u}|^2-2|\vec{u}||\vec{v}|\cos \angle(\vec{u}, \vec{v}) \\
	\vec{v}^{\,2}- 2\vec{u}\cdot \vec{v} + \vec{u}^{\,2} &= \vec{v}^{\,2}+\vec{u}^{\,2}-2|\vec{u}||\vec{v}|\cos \angle(\vec{u}, \vec{v}) \\
	\vec{u}\cdot \vec{v} &= |\vec{u}||\vec{v}|\cos \angle(\vec{u}, \vec{v})
}
\reg[Skalarproduktet II \label{scalardef2}]{
For to vektorer $ \vec{u} $ og $ \vec{v} $ er
\begin{equation*}
	\vec{u}\cdot \vec{v} = |\vec{u}||\vec{v}|\cos \angle(\vec{u}, \vec{v})
\end{equation*}
}
\eks{
For vektorene $\vec{u}$ og $\vec{v}$ har vi at $|\vec{u}|=4$, $|\vec{v}|=5$, og $\angle (\vec{u}, \vec{v})=60^\circ$. Finn skalarproduktet av $\vec{u}$ og $\vec{v}$.

\sv

Da $\cos 60^\circ=\frac{1}{2}$, har vi at
\[ \vec{u}\cdot\vec{v}=4\cdot5\cdot\frac{1}{2}=10 \]

}

\section{Vektorer vinkelrett på hverandre}
Fra \rref{scalardef2} kan vi gjøre en viktig observasjon; hvis $\angle (\vec{u},\vec{v})=90^\circ $, er $ {\cos\angle (\vec{u},\vec{v})=0} $, og da blir
\nn{
	\vec{u}\cdot\vec{v}=0
}
\reg[Vinkelrette vektorer \label{vinkrett} ]{
For to vektorer $ \vec{u} $ og $ \vec{v} $ har vi at
	\begin{equation*}
		\vec{u} \cdot \vec{v} = 0\iff \vec{u}\perp \vec{v}
	\end{equation*}
}
\spr{
Det er mange måter å uttrykke at $ {\vec{u}\perp \vec{v}} $ på. Blant annet kan vi si at
\begin{itemize}
	\item $ \vec{u} $ og $ \vec{v} $ står vinkelrette på hverandre.
	\item $ \vec{u} $ og $ \vec{v} $ står normalt på hverandre. 
	\item $\vec{u} $ er en normalvektor til $ \vec{v}$ (og omvendt).
	\item $\vec{u} $ og $ \vec{v}$ er ortogonale.
\end{itemize}
}
\eks[1]{
	Sjekk om vektorene $ {\vec{a}=[5, -3]} $, $ {\vec{b}=[6, 10]} $ og $ {\vec{c}=[2, 7] }$ er ortogonale.
	
	\sv 
	Vi har at
	\[ \vec{a}\cdot\vec{b}=5\cdot6+(-3)\cdot10=0 \]
	Altså er $ \vec{a}\perp\vec{b} $. Videre er
	\[ \vec{a}\cdot\vec{c}=5\cdot2+(-3)\cdot7 = 11 \]
	Altså er $ \vec{a} $ og $ \vec{c} $ \textsl{ikke} ortogonale. Videre er
	\[ \vec{b}\cdot\vec{c}=6\cdot2+10\cdot7=82 \]
	Altså er $\vec{b}$ og $\vec{c}$ \textsl{ikke} ortogonale.
}
\newpage
\info{Nullvektoren}{
I forkant av \rref{vinkrett} har vi bare argumentert for at\\ ${\vec{u}\perp\vec{v}\Rightarrow \vec{u}\cdot\vec{v}=0} $. For å rettferdiggjøre vilkåret som går begge veier, må vi vise at hvis ${ \vec{u}\cdot\vec{v} =0}$, så er $\vec{u}\perp\vec{v}$.\vsk

På intervallet $ [0^\circ, 180^\circ] $ er det bare ${\angle(\vec{u}, \vec{v})=90^\circ} $ som gir at $\cos \angle(\vec{u}, \vec{v})$. Skal skalarproduktet være 0 for andre vinkler, må derfor lengden av $ \vec{u} $ eller $ \vec{v} $ være 0. Den eneste vektoren med lengde 0 er \outl{nullvektoren} ${\vec{0}=[0, 0] }$, som rett og slett ikke har noen retning\footnote{Eventuelt kan man si at den peker i alle retninger!}. Det er likevel vanlig å definere at nullvektoren står vinkelrett på \textsl{alle} vektorer.
}
\section{Parallelle vektorer}
\regdef[Parallelle vektorer]{Hvis vinkelen mellom to vektorer er $ 0^\circ $ eller $ 180^\circ $, er de parallelle.}


\begin{figure}
\centering
\subfloat[]{\includegraphics{\figp{vekpara}}}\qquad
\subfloat[]{\includegraphics{\figp{vekparb}}}
\end{figure}
Gitt vektorene $ {\vec{u}=[x_1, y_1]} $ og $ {\vec{v}=[x_2, y_2]} $. La $ \theta $ og $ \alpha $ være vinkelen mellom $ x $-aksen og henholdsvis $ \vec{u} $ og $ \vec{v} $, med $ x $-aksen som høyre vinkelbein. Da er $ {\tan \theta =\frac{y_1}{x_1}} $ og $ {\tan \alpha =\frac{y_2}{x_2}} $. Hvis $ {\frac{y_1}{x_1}=\frac{y_2}{x_2}} $, er det to muligheter:
\begin{enumerate}[label=(\roman*)]
\item $ \theta=0^\circ $ og $ \alpha=180^\circ $, eller omvendt.
\item $ \theta = \alpha $
\end{enumerate}
I begge tilfeller er $ \angle(\vec{u}, \vec{v}) $ enten $ 0^\circ $ eller $ 180^\circ $, og da er $ \vec{u} $ og $ \vec{v} $ parallelle. Det omvendte gjelder også: Hvis punkt (i) eller (ii) gjelder, er $ {\frac{y_1}{x_1}=\frac{y_2}{x_2}} $. Det er ofte praktisk å omskrive denne sammehengen til forholdet mellom \outl{samsvarende komponenter}\footnote{
For vektorene $ [x_1, y_1] $ og $ [x_2, y_2] $ er disse samsvarende komponenter: \vspace{-5pt}
\begin{itemize}
	\item $ x_1 $ og $ x_2 $
	\item $ y_1 $ og $ y_2 $
\end{itemize}
}:
\begin{equation}\label{parforkl}
	\frac{x_1}{x_2}=\frac{y_1}{y_2}
\end{equation}
Det er også nyttig å merke seg at det må finnes to tall $ s $ og $ t $ slik at $ {\vec{u}=[sx_2, ty_2]} $. Hvis $ \vec{u}\parallel \vec{v}$, følger det av \eqref{parforkl} at $ {\frac{sx_2}{x_2}=\frac{ty_2}{y_2}}  $.
Dermed er $ s=t $. Omvendt; hvis $ \vec{u}=t[x_2, y_2] $, oppfyller $ \vec{u} $ og $ \vec{v} $ åpenbart \eqref{parforkl}.
\reg[Parallelle vektorer]{
	For to vektorer $ {\vec{u}=[x_1, y_1]} \text{ og } {\vec{v}=[x_2, y_2]} $ har vi at
	\begin{equation*}
		\frac{x_1}{x_2}=\frac{y_1}{y_2} \iff \vec{u}\parallel\vec{v}
	\end{equation*}
Alternativt, for et tall $ t $ har vi at
\begin{equation*}
 \vec{u}=t\vec{v}\iff \vec{u}\parallel\vec{v}
\end{equation*}
}
\spr{
Når $ {\vec{u}=t \vec{v}} $, sier vi at $ \vec{u} $ er et \outl{multiplum} av $ \vec{v} $ (og omvendt). Vi sier også at $ \vec{u} $ og $ \vec{v} $ er \outl{lineært avhengige}. \vsk

Om to vektorer ikke er parallelle, sier vi at de er \outl{lineært uavhengige}.
}
\eks[]{
Sjekk om $ {\vec{a}=[2, -3]} $ og $ {\vec{b}=[20, -45]} $ er parallelle med $ {\vec{c}=[10, -15]} $.

\sv
Vi har at
\nn{
\vec{c}=5[2, -3]=5\vec{a}
}
Dermed er $ \vec{a}\parallel \vec{c} $. Da $ \frac{20}{10}\neq \frac{-45}{-15} $, er $ \vec{b} $ og $ \vec{c} $ \textsl{ikke} parallelle.
}
\section{Vektorfunksjoner}
\subsection*{Parameterisering}
\regdef{
Gitt to funksjoner $ f(t) $ og $ g(t) $. En vektor $ \vec{v} $ på formen
\[ \vec{v}(t)=[f(t), g(t)] \]
er da en \outl{vektorfunksjon}.\vsk

Alternativt kan $ \vec{v} $ skrives som
\begin{equation*}\vec{v}(t): \left\lbrace{
		\begin{array}{l}
			x= f(t)   \\
			y= g(t)   
		\end{array}
	}\right. 
\end{equation*}
Mengden av alle endepunktene til $ \vec{v} $ utgjør grafen til $ \vec{v} $.
\fig{vek0}
}
\spr{
Koordinater uttrykt ved én eller flere variabler (parametre) kalles en \outl{parameterisering}\index{parameterisering}. En vektorfunksjoner er altså en parameterisering av en kurve.
}

\newpage
\subsection*{Vektorfunksjonen til en linje}
Gitt en linje beskrevet av vektorfunksjonen $ \vec{l}(t) $, som vist i figuren\\ under. Vi omtaler denne linja som linja $\vec{l}$.
\fig{linje}
Hvis en vektor $ \vec{r} $ er parallell med $ \vec{l} $, kalles den en \outl{retningsvektor}\index{retningsvektor}  for linja. Si at $ {\vec{r}=[a, b]} $ er en retningsvektor for $ \vec{l} $, og at $ A=(x_0, y_0) $ er et punkt på $ \vec{l} $. Om vi starter i $ A $ og vandrer parallellt med $ \vec{r} $, kan vi være sikre på at vi fortsatt befinner oss på linja. Dette må bety at vi for en variabel $ t $ kan nå et vilkårlig punkt $ {B=(x, y)} $ på linja ved følgende utregning:
\[B= A+t\vec{r} \]
På koordinatform kan vi skrive dette som
\[(x, y)= (x_0 +at, y_0+bt) \]
Altså kan linja skrives som en vektorfunksjon:\regv
\reg[Linje som vektorfunksjon]{
	En linje $ \vec{l}(t) $, som går gjennom punktet $ {A=(x_0, y_0,)} $ og har retningsvektor $ {\vec{r}=[a, b] }$, er gitt som
	\nn{
	\vec{l}=[x_0 + at, y_0+ bt]	\label{linje}
	}
}
\eks[1]{
En linje $\vec{l}(t)$ går gjennom punktene $(3, 6)$ og $(5, 13)$. Finn uttrykket til $\vec{l}$.

\sv
Vi finner retningsvektoren $\vec{r}$ til linja:
\[ \vec{r}=[5-3, 13-6]=[2, 7] \]
Altså er
\[ \vec{l}(t)=(3, 6)+[2, 7]t=[3+2t, 6+7t] \]
Alternativt kan vi skrive
\begin{equation*}\vec{l}(t): \left\lbrace{
		\begin{array}{l}
			x= 3+2t\\
			y= 6+7t   
		\end{array}
	}\right. 
\end{equation*}
\mers{Vi kunne like gjerne brukt $(5, 13)$ i stedet for $(3, 6)$ i uttrykket til $\vec{l}$.}
}
\eks[2]{
Gitt linja $\vec{l}(t)=[2+3t, 6-5t]$.
\abc{
\item Bruk uttrykket til $\vec{l}$ for å finne et punkt på linja, og for å finne en retningsvektor for linja.
\item Avgjør om punktet $(32, -60)$ ligger på linja.
}
\sv
\abc{
\item Av uttrykket til $\vec{l}$ ser vi at $(2, 6)$ er et punkt på linja, og at $[3, -5]$ er en retningsvektor for linja.
\item \metode{Metode 1}{2cm} \\
Skal punktet ligge på linja, må vektoren mellom $(2, 6)$ og $(32, -60)$ være en retningsvektor til linja. Denne vektoren er gitt som
\[ [32-2, -60-6]=[30, -66]=6[5, -11] \]
Denne vektoren er ikke parallell med $[3, -5]$, og dermed ligger ikke $(32, -60)$ på linja. \vsk

\metode{Metode 2}{2cm} \\
Skal $(32, -60)$ ligge på linja, må vi ha at
\begin{align*}
	2+3t&=32 \\
	t &= 10
\end{align*}
Men $6-5\cdot10\neq -60$, og dermed ligger ikke punktet på linja.
}
}
\section{Sirkellikningen}
Gitt en sirkel med sentrum $ S=(x_0, y_0) $, og et punkt $ A=(x, y) $. 
\fig{sirklig1}
Da er
\[ \vv{SA}=[x-x_0, y-y_0] \]
Av \rref{veklength} er da
\[ \left|\vv{SA}\right|^2=(x-x_0)^2+(y-y_0)^2 \]
Vi lar $r$ være radien til sirkelen. Hvis $A$ ligger på sirkelen, er $ \left|\vv{SA}\right|=r $, og dermed kan vi uttrykke $ r $ ved koordinatene til $ S $ og $ A $. Hvis $A$ ikke ligger på sirkelen, er $|\vv{SA}|\neq r$.\regv
\reg[Sirkellikningen]{
Gitt en sirkel med radius $ r $ og sentrum $ {S=(x_0, y_0)} $. Hvis punktet $ A=(x, y) $ ligger på buen til sirkelen, er
\nn{
(x-x_0)^2+(y-y_0)^2=r^2
}
Hvis likningen ikke er oppfylt, ligger ikke $A$ på sirkelbuen.
}
\newpage
\eks[]{
Finn sentrum og radien til sirkelen gitt av likningen
\begin{equation}\label{sirkligeks1}
	x^2 + y^2 - 4x + 10y - 20=0
\end{equation} \vs
\sv
Vi starter med å lage fullstendige kvadrat:
\alg{
x^2-4x &= (x-2)^2-2^2 \vn
y^2+10y &= (y+5)^2-5^2 
}
Altså kan vi skrive \eqref{sirkligeks1} som
\alg{
(x-2)^2+(y+5)^2-2^2-5^2-20&=0 \\
(x-2)^2+(y+5)^2&=7^2
}
Dermed har sirkelen sentrum $ (2, -5) $ og radius 7.
}

\section{Determinanter} \index{determinant}
\reg[\boldmath$ 2\times 2 $ determinanter]{
\outl{Determinanten} $ \det(\vec{u}, \vec{v}) $ av to vektorer $ {\vec{u}=[a, b] }$ og $ {\vec{v}=[c, d]} $ er gitt som
\begin{equation*}
	\det(\vec{u}, \vec{v}) = \left|\begin{matrix}
		a & b \\
		c & d
	\end{matrix}\right| = ad-bc
\end{equation*}
}
\eks{
	Gitt vektorene $ {\vec{u}=[-1, 3] }$ og $ {\vec{v}=[-2, 4]} $. Regn ut $ \det(\vec{u}, \vec{v})$.
	
	\sv
	\algv{
		\det(\vec{u}, \vec{v}) &= \left|\begin{matrix}
			-1 & 3 \\
			-2 & 4
		\end{matrix}\right| \\
		&= (-1)\cdot4-3\cdot(-2)	\\
		&= 2
	}
}
\reg[\detar \label{det2area}]{
Arealet $ A $ til et parallellogram formet av to vektorer $ \vec{u} $ og $ \vec{v} $ er gitt som
\begin{equation*}
	A= |\det(\vec{u}, \vec{v})|
\end{equation*}
\fig{utspprl}
Arealet $ A $ til en trekant formet av to vektorer $ \vec{u} $ og $ \vec{v} $ er gitt som
\begin{equation*}
	A= \frac{1}{2}|\det(\vec{u}, \vec{v})|
\end{equation*}
\fig{utsptrk}
}
\eks[]{
Finn arealet $A$ til parallellogrammet utspent av vektorene ${\vec{u}=[1, 1]}$ og ${\vec{v}=[-2, 5]}$.

\sv

Vi har at
\[ A=|\det(\vec{u}, \vec{v})|=|1\cdot5-1\cdot(-2)|=7 \]
}


\reg[\avstpunktlin \label{avstpunktlin}]{
	Avstanden $ h $ mellom et punkt $ B $ og en linje gitt av punktet $ A $ og retningsvektoren $ \vec{r} $ er gitt som
	\begin{equation*}
		h = \frac{\left|\det\left(\vv{AB}, \vec r\,\right)\right| }{|\vec{r}|}
	\end{equation*}
\fig{plin}
}
\newpage
\eks{
	Finn avstanden $h$ mellom punktet $B=(-2, 3)$ og linja \\
	$\vec{l}(t)=[1+t, -2+5t]$.
	
	\sv
	
	Linja er gitt av punktet $A=(1, -2)$ og retningsvektoren\\ $\vec{r}=[1, 5]$.
	Da $\vv{AB}=[-3, 5]$, har vi at
	\[ \det\left(\vv{AB}, \vec{r}\right)=(-3)\cdot5-5\cdot1=-20 \]
	Videre er $|\vec{r}|=\sqrt{1^2+5^2}=\sqrt{26}$. Følgelig har vi at
	\[ h= \frac{\left|\det\left(\vv{AB}, \vec{r}\right)\right|}{|\vec{r}|}=\frac{20}{\sqrt{26}}\]
}
\newpage
\fork{\ref{avstpunktlin} \avstpunktlin}{
La en linje $ \vec{l}(t) $ i rommet være gitt  av et punkt $ A $ og en retningsvektor $ \vec{r} $. I tillegg ligger et punkt \textit{B} utenfor linja, som vist i figuren under
\fig{plin}
Den korteste avstanden fra \textit{B} til linja er høyden \textit{h} i trekanten formet av $ \vec{r} $ og $ \vv{AB} $. Arealet til denne trekanten er gitt av \rref{det2area}:
\[ \frac{1}{2}\left|\det\left(\vv{AB}, \vec r\right)\right| \]
Av den klassiske arealformelen for en trekant (se \mb) har vi at
\alg{\frac{1}{2}|\vec{r}\,|h &=\frac{1}{2}\left|\det\left(\vv{AB}, \vec r\right)\right|  \br
	h &= \frac{\left|\det\left(\vv{AB}, \vec r\right)\right| }{|\vec{r}|}
}
}
\section{Forklaringer}
\fork{\ref{det2area} \detar }{
Vi lar $A_N$ betegne arealet til en geometrisk form $N$.
\begin{figure}
\subfloat[]{\includegraphics{\figp{utspprl}}}\qquad
\subfloat[]{\includegraphics{\figp{utspprlforkl}}}\qquad
\end{figure}
Gitt to vektorer $ {\vec{u}=[a, b]} $ og $ {\vec{v}=[c, d]} $, hvor $ {a, b, c, d >0} $, som vist i figur \textsl{(a)}. Plasserer vi vektorene i grunnstillingen er punktene vist i figur \textsl{(b)} gitt som
\alg{
O &= (0, 0) & B&=(a, b) & C &= (a+b, c+d) \br
  D&= (c, d) & E &= (a+c, 0) & F &= (0, b+d)
}
Med $ OE $ som grunnlinje, har $ \triangle OEB $ høyde $ b $, altså er
\[  2A_{\triangle OEB}=(a+c)b \]
Tilsvarende er
\[  2A_{\triangle FDO}=(b+d)c \]
Da $ A_{\triangle OEB}=A_{\triangle CDF}$ og $ A_{\triangle FDO}=A_{\triangle EBC} $, har vi at
\alg{
A_{\square ABCD}&=A_{\square OECF}-2A_{\triangle OBE}-2A_{\triangle FDO} \\
&=(a+c)(b+d)-(a+c)b-(b+d)c \\
&=(a+c)d-(b+d)c\\
&= ad-bc
}
I figurene har vi antatt at (den minste) vinkelen mellom $ \vec{v} $ og $ x $-aksen er mindre enn vinkelen mellom $ \vec{u} $ og $ x $-aksen. I omvendt tilfelle ville vi fått at 
\[ A_{\square OECF}=bc-ad \]
Altså er
\[  A_{\square OECF}=|ad-bc| \]
På lignende måte kan det vises at dette gjelder for alle $a, b, c, d\in\mathbb{R} $. \vsk

\mers{Dette kan også vises veldig kortfattet ved å anvende trigonometri. Dette er gitt som en oppgave i \tmto.}
}
\end{document}